%Planning to understand/simplify/improve. If you are interested, please contact me (Evangelos Kipouridis), I would be more than happy to do that collaboratively.
%Licensed under CC BY 4.0
\documentclass[11pt]{article}

\usepackage[margin=1in]{geometry}
\usepackage{amsmath,amssymb,amsthm,mathtools}
\usepackage{microtype}
\usepackage[hidelinks]{hyperref}

\newtheorem{theorem}{Theorem}[section]
\newtheorem{lemma}[theorem]{Lemma}
\newtheorem{proposition}[theorem]{Proposition}
\newtheorem{corollary}[theorem]{Corollary}
\theoremstyle{definition}
\newtheorem{definition}[theorem]{Definition}
\newtheorem{remark}[theorem]{Remark}

\newcommand{\R}{\mathbb{R}}
\newcommand{\Z}{\mathbb{Z}}
\newcommand{\N}{\mathbb{N}}
\newcommand{\cU}{\mathcal{U}}
\newcommand{\cH}{\mathcal{H}}
\newcommand{\OPT}{\operatorname{OPT}}
\newcommand{\cost}{\operatorname{cost}}
\newcommand{\supp}{\operatorname{supp}}
\newcommand{\lca}{\operatorname{lca}}
\newcommand{\Est}{\operatorname{Est}}
\newcommand{\poly}{\operatorname{poly}}
\newcommand{\polylog}{\operatorname{polylog}}
\newcommand{\Fp}{F_p}

\title{One-Pass Approximation of Best-Fit Ultrametrics\\
with Unrestricted Decoding Time}
\author{}
\date{}

\begin{document}

\maketitle

\begin{abstract}
Let \(V\) be a set of \(n\) labels, let \(E=\binom{V}{2}\), and let
\(D:E\to\{1,\ldots,2^w-1\}\) be given by a turnstile stream of coordinate
updates.  We study the problem of finding an ultrametric \(U\) on \(V\) that
minimizes \(\lVert D-U\rVert_q\).  For each \(q\in\{0,1,2\}\), we give a
randomized one-pass algorithm that returns a \((1+\varepsilon)\)-approximate
ultrametric.  If \(w=O(\log n)\), the number of stream updates is polynomial in
\(n\), and the required failure probability is \(n^{-O(1)}\), then the
algorithm uses \(\widetilde O(n/\varepsilon^2)\) bits and
\(n^{O(n)}(1/\varepsilon)^{O(1)}\) decoding time.  For fixed
\(\varepsilon\), this decoding time is \(n^{O(n)}\).  The algorithm is based
on a finite family of
canonical ultrametrics and a linear norm sketch that is simultaneously correct
for every member of that family.

We also prove a lower bound for reporting the optimum value.  For every fixed
real number \(p>4\), a randomized one-pass strict-turnstile algorithm that
returns a sufficiently accurate multiplicative estimate of
\(\min_U\lVert D-U\rVert_p\) requires
\(\widetilde\Omega(n^{2-4/p})\) bits.  The post-processing time of the
algorithm is unrestricted in this lower bound.
\end{abstract}

\section{Introduction}

An input consists of one number \(D(\{u,v\})\) for every unordered pair of
distinct labels \(u,v\in V\).  The output is an ultrametric \(U\) on \(V\).
The error of \(U\) is measured by applying an \(\ell_q\) objective to the
vector \(D-U\), indexed by the pairs in \(\binom{V}{2}\).  The input vector is
not stored explicitly.  Instead, its coordinates are changed by a sequence of
updates, and the algorithm reads that sequence once.

The following theorem is the positive result of this paper.  All terms in its
statement are defined in Section~\ref{sec:preliminaries}.

\begin{theorem}[Upper bound]\label{thm:upper-intro}
Fix \(q\in\{0,1,2\}\), \(0<\varepsilon<1\), and \(0<\delta<1/2\).
There is a randomized one-pass turnstile algorithm that, with probability at
least \(1-\delta\), returns an ultrametric \(\widehat U\) satisfying
\[
  \lVert D-\widehat U\rVert_q
  \le (1+\varepsilon)\min_{U\in\cU(V)}\lVert D-U\rVert_q.
\]
The algorithm uses
\[
 O\!\left(
   \varepsilon^{-2}
   \bigl(n(\log n+w)+\log(1/\delta)\bigr)
   \bigl(\log n+w+\log(1/\varepsilon)\bigr)^{O(1)}
 \right)
\]
bits.  Its decoding time is
\[
 2^{O(n(\log n+w))}
 \bigl(n+1/\varepsilon+\log(1/\delta)\bigr)^{O(1)}.
\]
In particular, if \(w=O(\log n)\) and
\(\delta=n^{-O(1)}\), then the space is
\(\widetilde O(n/\varepsilon^2)\) bits and the decoding time is
\(n^{O(n)}(1/\varepsilon)^{O(1)}\).  For fixed \(\varepsilon\), the decoding
time is \(n^{O(n)}\).
\end{theorem}

Here and below, \(\cU(V)\) denotes the set of ultrametrics on \(V\).  Theorem
\ref{thm:upper-intro} allows arbitrary negative intermediate coordinates in
the stream, subject to the requirement that the final vector is a positive
\(w\)-bit distance array.  It therefore also applies to strict-turnstile and
insertion-only streams.

The following theorem is the negative result of this paper.  It concerns an
algorithm that reports the optimum value; it does not concern an algorithm
that returns only an ultrametric without reporting its error.

\begin{theorem}[Optimum-value lower bound]\label{thm:lower-intro}
Fix a real number \(p>4\).  There is a constant
\(\varepsilon_p>0\), depending only on \(p\), such that the following holds.
Every randomized one-pass strict-turnstile algorithm that, with probability at
least \(2/3\), returns a number \(Z\) satisfying
\[
 (1-\varepsilon_p)\OPT_p(D)
 \le Z\le
 (1+\varepsilon_p)\OPT_p(D),
 \qquad
 \OPT_p(D):=\min_{U\in\cU(V)}\lVert D-U\rVert_p,
\]
uses \(\widetilde\Omega(n^{2-4/p})\) bits on a family of streams whose
lengths and update magnitudes are polynomial in \(n\).  The lower bound holds
without a restriction on the post-processing time.  Every final distance in
the constructed family has \(O(\log n)\) bits.
\end{theorem}

For \(p>4\), the exponent \(2-4/p\) is greater than \(1\).  Hence Theorem
\ref{thm:lower-intro} excludes semi-streaming space for the optimum-value
problem.  The proof uses the standard lower bound for estimating a high
frequency moment and an exact identity between that frequency moment and the
optimum fitting cost of the constructed distance array.

Ultrametric fitting in the semi-streaming model was introduced by Carmel,
Das, Kipouridis, and Pipis~\cite{CarmelDasKipouridisPipis2025}.  The norm
sketches used here are standard turnstile sketches; see
Cormode, Datar, Indyk, and Muthukrishnan~\cite{CormodeDatarIndykMuthukrishnan2003},
Indyk~\cite{Indyk2006}, and Kane, Nelson, and
Woodruff~\cite{KaneNelsonWoodruff2010Norms,KaneNelsonWoodruff2010Distinct}.

\section{Definitions and notation}\label{sec:preliminaries}

\subsection{Finite sets and vectors}

For a positive integer \(r\), write
\([r]:=\{1,\ldots,r\}\).  Let \(V\) be a finite set and let
\(n:=|V|\ge 2\).  Define
\[
 E:=\binom{V}{2}
   :=\bigl\{\{u,v\}:u,v\in V,\ u\ne v\bigr\},
 \qquad
 N:=|E|=\binom n2.
\]
A vector indexed by \(E\) is identified with a function from \(E\) to
\(\R\).  For \(x\in\R^E\), define
\[
 \lVert x\rVert_0
 :=\bigl|\{e\in E:x(e)\ne 0\}\bigr|,
 \qquad
 \lVert x\rVert_q
 :=\left(\sum_{e\in E}|x(e)|^q\right)^{1/q}
 \quad(q\ge 1).
\]
Thus \(\lVert x\rVert_0\) is the number of nonzero coordinates of \(x\).
The notation \(\lVert\cdot\rVert_0\) does not denote a norm.

All logarithms have base \(2\).  The notation \(f(n)=\widetilde O(g(n))\)
means that there are constants \(c,C>0\) such that
\(f(n)\le Cg(n)(\log n)^c\) for all sufficiently large \(n\).  The notation
\(f(n)=\widetilde\Omega(g(n))\) means that there are constants \(c,C>0\)
such that \(f(n)\ge Cg(n)/(\log n)^c\) for all sufficiently large \(n\) in
the stated family of input sizes.
For positive real numbers \(z_1,\ldots,z_k\), the notation
\(\polylog(z_1,\ldots,z_k)\) denotes
\[
 \left(1+\sum_{i=1}^k\log(1+z_i)\right)^c
\]
for an absolute constant \(c>0\).  The value of \(c\) may differ between
occurrences of the notation.

\subsection{Ultrametrics and dendrograms}

\begin{definition}[Ultrametric]
An ultrametric on \(V\) is a function \(U:E\to\R_{\ge 0}\) such that, for
every three distinct labels \(a,b,c\in V\),
\[
 U(\{a,c\})
 \le
 \max\bigl\{U(\{a,b\}),U(\{b,c\})\bigr\}.
\]
The set of ultrametrics on \(V\) is denoted by \(\cU(V)\).
\end{definition}

The value zero is allowed in this definition.  Allowing zero does not change
the optimum for the positive input arrays considered in this paper.

\begin{definition}[Rooted tree terminology]
A rooted tree is a finite connected acyclic graph with a distinguished vertex,
called its root.  If \(x\) and \(y\) are vertices of a rooted tree, then \(x\)
is an ancestor of \(y\) if \(x\) lies on the unique path from the root to
\(y\).  A child of \(x\) is a neighbor \(y\) of \(x\) for which \(x\) is the
ancestor of \(y\).  A leaf is a vertex with no child.  An internal vertex is a
vertex that is not a leaf.  The least common ancestor of two vertices \(u,v\)
is the common ancestor of \(u\) and \(v\) that has no strict descendant that
is also a common ancestor; it is denoted by \(\lca_T(u,v)\).
\end{definition}

\begin{definition}[Dendrogram]
A dendrogram on \(V\) is a pair \((T,h)\) with the following properties.
The object \(T\) is a rooted tree whose leaves are in bijection with \(V\),
every internal vertex of \(T\) has at least two children, and
\(h:I(T)\to\R_{\ge0}\) is a function on the set \(I(T)\) of internal vertices
such that
\[
 h(x)\ge h(y)
 \quad\text{whenever \(x\) is an ancestor of \(y\).}
\]
The ultrametric represented by \((T,h)\) is the function
\(U_{T,h}:E\to\R_{\ge0}\) defined by
\[
 U_{T,h}(\{u,v\}):=h(\lca_T(u,v)).
\]
\end{definition}

\begin{lemma}[Dendrogram representation]\label{lem:dendrogram-representation}
Every dendrogram on \(V\) represents an ultrametric on \(V\).  Every
ultrametric on \(V\) is represented by a dendrogram with at most \(n-1\)
internal vertices.
\end{lemma}

\begin{proof}
Let \((T,h)\) be a dendrogram and let \(a,b,c\in V\) be distinct.  Among the
three least common ancestors of the three pairs, one is an ancestor of the
other two.  The height of that ancestor is the maximum of the three pair
heights and occurs for at least two of the three pairs.  Therefore every one
of the three pair heights is at most the maximum of the other two.  This is
the ultrametric inequality.

Conversely, let \(U\) be an ultrametric.  For \(t\ge0\), define
\[
 u\sim_t v
 \quad\Longleftrightarrow\quad
 u=v\text{ or }U(\{u,v\})\le t.
\]
The ultrametric inequality implies that \(\sim_t\) is transitive.  It is also
reflexive and symmetric, so it is an equivalence relation.  If \(s\le t\),
then every equivalence class of \(\sim_s\) is contained in an equivalence
class of \(\sim_t\).  Take all distinct equivalence classes obtained when
\(t\) ranges over \(\{U(e):e\in E\}\), and add every singleton set.  Ordered
by inclusion, these sets form a laminar family containing \(V\).  Its Hasse
diagram, directed from \(V\) toward the singleton sets, is a rooted tree with
leaf set \(V\).  Give each non-singleton set \(C\) the smallest value \(t\)
for which \(C\) is an equivalence class of \(\sim_t\).  These values are
nonincreasing from an ancestor to a descendant.  For distinct \(u,v\in V\),
the smallest set in the family containing both labels has height
\(U(\{u,v\})\).  Hence the resulting dendrogram represents \(U\).

Each internal vertex has at least two children.  A rooted tree with \(n\)
leaves and at least two children at every internal vertex has at most \(n-1\)
internal vertices.
\end{proof}

\subsection{The fitting objectives}

Fix a positive integer \(w\) and define \(B:=2^w-1\).  A distance array on
\(V\) is a function
\[
 D:E\to[B].
\]
No triangle inequality is imposed on \(D\).  For \(q\in\{0,1,2\}\) and
\(U\in\cU(V)\), define
\[
 \cost_q(D,U):=\lVert D-U\rVert_q,
 \qquad
 \OPT_q(D):=\min_{U\in\cU(V)}\cost_q(D,U).
\]
For a real number \(p\ge1\), the same notation is used with the
\(\ell_p\) norm:
\[
 \cost_p(D,U):=
 \left(\sum_{e\in E}|D(e)-U(e)|^p\right)^{1/p},
 \qquad
 \OPT_p(D):=\min_{U\in\cU(V)}\cost_p(D,U).
\]

\subsection{Streaming models}

\begin{definition}[Turnstile stream]
A turnstile stream for a distance array is a finite sequence
\[
 (e_1,\Delta_1),\ldots,(e_L,\Delta_L)
 \quad\text{with}\quad
 e_t\in E,
 \quad
 \Delta_t\in\Z.
\]
The vector after update \(t\) is \(D^{(t)}\in\Z^E\), defined by
\[
 D^{(0)}(e):=0,
 \qquad
 D^{(t)}(e):=D^{(t-1)}(e)+
 \begin{cases}
 \Delta_t,&e=e_t,\\
 0,&e\ne e_t.
 \end{cases}
\]
The final vector is \(D:=D^{(L)}\).  A strict-turnstile stream is a turnstile
stream satisfying \(D^{(t)}(e)\ge0\) for every \(t\in[L]\) and \(e\in E\).
An insertion-only stream is a strict-turnstile stream satisfying
\(\Delta_t\ge0\) for every \(t\in[L]\).
\end{definition}

The upper bound assumes that \(D(e)\in[B]\) for every \(e\in E\), that
\(L\le n^{O(1)}\), and that \(|\Delta_t|\le 2^{O(w)}\).  The intermediate
coordinates of a turnstile stream may be negative.  The lower bound uses only
strict-turnstile streams.

A one-pass streaming algorithm reads the updates in their given order and
does not read an update a second time.  Its space is the maximum number of bits
in its internal state during the stream and its post-processing.  Its decoding
time is the number of elementary arithmetic and comparison operations used
after the final update.  A semi-streaming algorithm uses
\(O(n\polylog n)\) bits.  The input stream is fixed independently of the
private random bits used by the algorithm.
An integer parameter is polynomially bounded in \(n\) if it is at most
\(n^c\) for a constant \(c>0\).

\section{Linear norm sketches}\label{sec:sketches}

The proof of the upper bound uses a standard streaming primitive.  We state
the primitive in the form required here.

A randomized linear sketch of a vector \(x\) consists of a random linear map
\(S\) and the value \(Sx\).  Linearity means
\(S(ax+by)=aSx+bSy\) whenever the displayed operations are defined.

\begin{lemma}[Linear sketches for \(q\in\{0,1,2\}\)]
\label{lem:norm-sketch}
Let \(d,M,L\) be positive integers, let \(q\in\{0,1,2\}\), and let
\(0<\eta<1/2\) and \(0<\rho<1/2\).  There is a randomized linear sketch
\(S_q\) and a decoder \(\Est_q\) with the following properties.

For each fixed vector \(x\in\mathbb{Q}^d\) whose coordinates have numerators
and denominators of absolute value at most \(M\),
\[
 \Pr\left[
  (1-\eta)\lVert x\rVert_q
  \le \Est_q(S_qx)
  \le (1+\eta)\lVert x\rVert_q
 \right]
 \ge 1-\rho.
\]
For \(q=0\), the expression \(\lVert x\rVert_0\) has the meaning given in
Section~\ref{sec:preliminaries}.  If \(x=0\), then
\(\Est_q(S_qx)=0\) with probability \(1\).

The sketch is linear: whenever \(x,y\in\mathbb{Q}^d\) and
\(a,b\in\mathbb{Q}\), the sketch of \(ax+by\) can be computed from the
sketches of \(x\) and \(y\).  A stream of \(L\) coordinate updates to an
integer vector can be processed in one pass.  The number of stored bits is
\[
 O\!\left(
  \eta^{-2}\log(1/\rho)
  \polylog(d,M,L,1/\eta)
 \right).
\]
Given a rational query vector \(y\) of the stated coordinate size, \(S_qy\)
can be computed from the random seed describing \(S_q\), without access to the
stream.
\end{lemma}

\begin{proof}
For \(q=2\), one may use the sketch of Alon, Matias, and
Szegedy~\cite{AlonMatiasSzegedy1999}.  One row stores the linear measurement
\[
 Z:=\sum_{i=1}^d\sigma_i x_i,
\]
and the decoder uses \(Z^2\), where the signs
\(\sigma_i\in\{-1,1\}\) are uniform and four-wise independent, meaning that
every subfamily of at most four signs is mutually independent.  The
expectation of \(Z^2\) is \(\lVert x\rVert_2^2\), and its variance is at most
\(2\lVert x\rVert_2^4\).  Averaging \(O(\eta^{-2})\) independent copies gives
constant failure probability by Chebyshev's inequality.  Taking the median of
\(O(\log(1/\rho))\) independent averages gives failure probability at most
\(\rho\).  Taking a square root gives the stated estimate of
\(\lVert x\rVert_2\), after changing the constant in the accuracy parameter.

For \(q=1\), one may use independent Cauchy projections as in
Indyk~\cite{Indyk2006}.  If \(C_1,\ldots,C_d\) are independent standard
Cauchy random variables, where a standard Cauchy random variable has density
\[
 f(t):=\frac{1}{\pi(1+t^2)}
 \qquad(t\in\R),
\]
then
\[
 \sum_{i=1}^d x_iC_i
 \quad\text{and}\quad
 \lVert x\rVert_1 C_1
\]
have the same distribution.  The median of \(|C_1|\) is \(1\), and the
density of \(|C_1|\) is positive and continuous in a neighborhood of \(1\).
The Chernoff bound for the empirical median therefore gives relative error
\(\eta\) from \(O(\eta^{-2}\log(1/\rho))\) independent projections.

For \(q=0\), use a turnstile Hamming-norm sketch; such sketches are given by
Cormode, Datar, Indyk, and
Muthukrishnan~\cite{CormodeDatarIndykMuthukrishnan2003} and by Kane, Nelson,
and Woodruff~\cite{KaneNelsonWoodruff2010Distinct}.  These sketches consist of
linear measurements and estimate the number of nonzero coordinates.  If the
coordinates are rational, choose a prime larger than every possible nonzero
numerator and every denominator, and interpret every coordinate in the
corresponding finite field.  This preserves whether each coordinate is zero.

In all three cases, each measurement is a linear function of \(x\).  Hence a
coordinate update changes the stored measurements without storing \(x\).
Independent repetition gives the factor \(\log(1/\rho)\).  Standard bounded
independence and finite-precision implementations use the stated number of
bits for vectors and streams with the stated magnitude bounds; see
\cite{Indyk2006,KaneNelsonWoodruff2010Norms,KaneNelsonWoodruff2010Distinct}.
\end{proof}

\section{Canonical ultrametrics}\label{sec:canonical}

This section constructs a finite family that contains an optimum solution for
each of the three objectives.

Fix a rooted tree \(T\) whose leaves are \(V\) and whose internal vertices
have at least two children.  For each internal vertex \(v\in I(T)\), define
\[
 P_v:=\bigl\{\{a,b\}\in E:\lca_T(a,b)=v\bigr\}.
\]
The sets \(P_v\), for \(v\in I(T)\), form a partition of \(E\), and each
\(P_v\) is nonempty.

\begin{lemma}[Existence of an optimum on a fixed tree]
\label{lem:existence-fixed-tree}
For \(q=0\) and for every real \(q\ge1\), every distance array
\(D:E\to[B]\), and every rooted tree \(T\) considered above, there is a
feasible height function \(h\) that minimizes
\(\cost_q(D,U_{T,h})\).
\end{lemma}

\begin{proof}
For \(q=0\), every feasible height function has a cost in the finite set
\(\{0,\ldots,N\}\).  The nonempty set of attained costs therefore has a
smallest member.

For every real \(q\ge1\), the feasible height functions form a closed subset
of \(\R_{\ge0}^{I(T)}\), and the objective is a continuous function of the
heights.  Because every \(P_v\) is nonempty, the objective tends to infinity
when the largest height tends to infinity.
Every nonempty sublevel set is therefore closed and bounded, hence compact.
A continuous function on a nonempty compact set attains its minimum.
\end{proof}

\begin{lemma}[Canonical heights]\label{lem:canonical-heights}
Let \(D:E\to[B]\), and fix a rooted tree \(T\) as above.

\begin{enumerate}
\item There is a height function minimizing
\(\cost_0(D,U_{T,h})\) whose values belong to
\[
 \cH_0(D):=\{D(e):e\in E\}.
\]

\item There is a height function minimizing
\(\cost_1(D,U_{T,h})\) whose values belong to
\[
 \cH_1(D):=\{D(e):e\in E\}.
\]

\item There is a height function minimizing
\(\cost_2(D,U_{T,h})\) such that every height is the arithmetic mean of a
nonempty multiset of entries of \(D\).  In particular, every height belongs to
\[
 \cH_2(B,N):=
 \left\{\frac{s}{k}:
  1\le k\le N,
  k\le s\le kB,
  s\in\Z
 \right\}.
\]
\end{enumerate}
\end{lemma}

\begin{proof}
For the first claim, let \(h\) be any feasible height function.  Define
\(\phi:\R_{\ge0}\to\cH_0(D)\) by
\[
 \phi(t):=
 \begin{cases}
  \min\cH_0(D),&t<\min\cH_0(D),\\
  \max\{s\in\cH_0(D):s\le t\},&t\ge\min\cH_0(D).
 \end{cases}
\]
The function \(\phi\) is nondecreasing and satisfies \(\phi(s)=s\) for every
\(s\in\cH_0(D)\).  Therefore \(\phi\circ h\) is feasible.  If
\(U_{T,h}(e)=D(e)\), then
\(U_{T,\phi\circ h}(e)=D(e)\).  Thus every agreement is preserved, and
\[
 \cost_0(D,U_{T,\phi\circ h})
 \le \cost_0(D,U_{T,h}).
\]
Applying this transformation to an optimum \(h\) proves the first claim.

For the remaining claims, connect two internal vertices by an edge when one
is the parent of the other.  For a feasible height function \(h\), an equality
block is a connected component of the graph containing precisely those
parent--child edges \(xy\) for which \(h(x)=h(y)\).  All vertices in an
equality block \(C\) have one common height, denoted by \(t_C\).  Define
\[
 P_C:=\bigcup_{v\in C}P_v.
\]
Because adjacent equality blocks have different heights, the value \(t_C\)
can be increased or decreased by a sufficiently small positive amount while
all other block heights remain fixed, unless \(t_C=0\).  The case \(t_C=0\)
cannot occur in an optimum for the \(\ell_1\) or \(\ell_2\) objective, because
all entries of \(D\) are positive and a sufficiently small increase would
strictly decrease the objective.

For the \(\ell_2\) objective, minimizing the norm is equivalent to minimizing
its square.  The terms depending on \(t_C\) are
\[
 f_C(t):=\sum_{e\in P_C}(D(e)-t)^2.
\]
Since \(t_C\) is an interior local minimizer,
\[
 0=f_C'(t_C)=2\sum_{e\in P_C}(t_C-D(e)).
\]
Consequently,
\[
 t_C=\frac{1}{|P_C|}\sum_{e\in P_C}D(e).
\]
The denominator is between \(1\) and \(N\), and the numerator is an integer
between \(|P_C|\) and \(|P_C|B\).  This proves the third claim.

For the \(\ell_1\) objective, choose among all optimum height functions one
with the smallest number of equality blocks; subject to this choice, choose
one with the largest number of canonical equality blocks, where an equality
block \(C\) is called canonical when \(t_C\in\cH_1(D)\).  These choices are
possible by Lemma~\ref{lem:existence-fixed-tree}.  Suppose that a block \(C\)
has height
\(t_C\notin\cH_1(D)\).  The terms depending on \(t_C\) are
\[
 g_C(t):=\sum_{e\in P_C}|D(e)-t|.
\]
The function \(g_C\) is linear on every open interval containing no value
from \(\cH_1(D)\).  Since \(t_C\) is an interior local minimizer, the slope on
the interval containing \(t_C\) is zero.  Choose a direction from \(t_C\)
toward a closest value of \(\cH_1(D)\).  Move \(t_C\) in this direction,
without changing any other block height, until either \(t_C\) reaches that
value or \(t_C\) reaches the height of an adjacent equality block.
The objective does not change during this movement.  In the first case, the
number of canonical blocks increases.  In the second case, two equality
blocks merge.  Both conclusions contradict the choice of \(h\).  Hence every
block height belongs to \(\cH_1(D)\), which proves the second claim.
\end{proof}

The families in Lemma~\ref{lem:canonical-heights} that depend on \(D\) are
contained in the universal set \([B]\).  Define
\[
 \cH_0:=\cH_1:=[B],
 \qquad
 \cH_2:=\cH_2(B,N).
\]
For \(q\in\{0,1,2\}\), let \(\mathcal C_q(V,B)\) be the set of ultrametrics
that have a dendrogram representation with at most \(n-1\) internal vertices
and all internal heights in \(\cH_q\).

\begin{lemma}[Number of candidates]\label{lem:candidate-count}
For every \(q\in\{0,1,2\}\),
\[
 |\mathcal C_q(V,B)|
 \le 2^{O(n(\log n+w))}.
\]
Moreover, \(\mathcal C_q(V,B)\) contains an ultrametric attaining
\(\OPT_q(D)\).
\end{lemma}

\begin{proof}
Every rooted tree under consideration has \(n\) labelled leaves and at most
\(n-1\) internal vertices.  After assigning distinct auxiliary labels to the
internal vertices, such a tree can be described by at most \(2n-2\) parent
pointers, each taking at most \(2n-1\) values.  Therefore the number of rooted
tree topologies is at most
\[
 (2n)^{2n}=2^{O(n\log n)}.
\]
The auxiliary labels can only overcount the number of trees.

For \(q\in\{0,1\}\), \( |\cH_q|=B<2^w\).  For \(q=2\),
\[
 |\cH_2|
 \le\sum_{k=1}^N kB
 \le BN^2.
\]
For each topology, at most \(n-1\) heights are selected.  Ignoring the
monotonicity constraint can only increase the count, so
\[
 |\mathcal C_q(V,B)|
 \le (2n)^{2n}|\cH_q|^{n-1}
 \le 2^{O(n(\log n+w))}.
\]
The existence of an optimum member follows from the dendrogram
representation of ultrametrics and Lemma~\ref{lem:canonical-heights}.
\end{proof}

\section{One-pass upper bounds}\label{sec:upper}

We now prove Theorem~\ref{thm:upper-intro}.

\begin{proof}[Proof of Theorem~\ref{thm:upper-intro}]
Fix \(q\in\{0,1,2\}\) and let
\[
 K:=|\mathcal C_q(V,B)|.
\]
Set
\[
 \eta:=\frac{\varepsilon}{2+\varepsilon}
 \quad\text{and}\quad
 \rho:=\frac{\delta}{K}.
\]
Choose the sketch \(S_q\) of Lemma~\ref{lem:norm-sketch} with parameters
\(\eta,\rho\).  During the stream, maintain only \(S_qD^{(t)}\).  Linearity
allows every update to be applied directly to this state.

The sketch dimension is \(d=N\).  For \(q\in\{0,1\}\), every candidate
coordinate is an integer in \([B]\).  For \(q=2\), every candidate coordinate
has a numerator of absolute value at most \(NB\) and a denominator at most
\(N\).  Thus every coordinate of \(D-U\) has a numerator and denominator of
absolute value at most \(2NB\), and Lemma~\ref{lem:norm-sketch} applies with
\(M=2NB\).  In particular, each candidate sketch can be evaluated using the
bit complexity included in that lemma.

After the final update, enumerate all candidates
\(U\in\mathcal C_q(V,B)\).  For each candidate, compute \(S_qU\) from the
description of \(U\) and the random seed describing \(S_q\), and define
\[
 \widehat c(U):=\Est_q(S_qD-S_qU)
               =\Est_q(S_q(D-U)).
\]
Return a candidate \(\widehat U\) minimizing \(\widehat c(U)\).  Ties are
resolved according to any fixed ordering of the finite candidate family.

For each fixed \(U\in\mathcal C_q(V,B)\), Lemma
\ref{lem:norm-sketch} gives
\[
 \Pr\left[
 (1-\eta)\cost_q(D,U)
 \le\widehat c(U)\le
 (1+\eta)\cost_q(D,U)
 \right]\ge1-\rho.
\]
The candidate family is fixed independently of the random sketch.  The union
bound therefore implies that, with probability at least
\(1-K\rho=1-\delta\), the displayed inequalities hold simultaneously for
every candidate.

Condition on this simultaneous event.  By Lemma
\ref{lem:candidate-count}, there is a candidate \(U^\star\) satisfying
\(\cost_q(D,U^\star)=\OPT_q(D)\).  Hence
\begin{align*}
 (1-\eta)\cost_q(D,\widehat U)
 &\le \widehat c(\widehat U)\\
 &\le \widehat c(U^\star)\\
 &\le (1+\eta)\cost_q(D,U^\star).
\end{align*}
Since
\[
 \frac{1+\eta}{1-\eta}=1+\varepsilon,
\]
the required approximation follows.  If \(\OPT_q(D)=0\), then
\(S_q(D-U^\star)=0\), so its estimate is exactly zero.  Simultaneous
correctness then implies \(\cost_q(D,\widehat U)=0\).

By Lemma~\ref{lem:candidate-count},
\[
 \log(1/\rho)=\log K+\log(1/\delta)
 =O\bigl(n(\log n+w)+\log(1/\delta)\bigr).
\]
Substitution into Lemma~\ref{lem:norm-sketch} proves the stated space bound.
Every candidate has at most \(n-1\) internal heights, each represented using
\(O(w+\log n)\) bits.  Its vector of pair values and its sketch can be
computed in polynomial time.  Enumerating the candidate family therefore
takes
\[
 2^{O(n(\log n+w))}
 \bigl(n+1/\varepsilon+\log(1/\delta)\bigr)^{O(1)}
\]
time and does not require the candidates to be stored simultaneously.  This
proves the theorem.
\end{proof}

\section{A strict-turnstile lower bound for reporting the optimum value}
\label{sec:lower}

We reduce from frequency-moment estimation.  For a positive integer \(d\), a
frequency vector is a vector \(x\in\Z_{\ge0}^d\).  It is presented by an
insertion-only stream whose update \(i\in[d]\) replaces \(x_i\) by
\(x_i+1\).  For a fixed real number \(p>2\), define
\[
 \Fp(x):=\sum_{i=1}^d x_i^p.
\]
We use the following standard lower bound; see Bar-Yossef, Jayram, Kumar, and
Sivakumar~\cite{BarYossefJayramKumarSivakumar2004} and
Ganguly~\cite{Ganguly2011}.

\begin{theorem}[High-frequency-moment lower bound]
\label{thm:fp-lower}
Fix a real number \(p>2\).  There is a constant \(\gamma_p>0\) such that every
randomized one-pass algorithm that, with probability at least \(2/3\), returns
a number \(Y\) satisfying
\[
 (1-\gamma_p)\Fp(x)\le Y\le(1+\gamma_p)\Fp(x)
\]
uses \(\widetilde\Omega(d^{1-2/p})\) bits.  The lower bound holds for
insertion-only streams of polynomial length, and the resulting coordinates
of \(x\) are bounded by a polynomial in \(d\).
\end{theorem}

We next define the distance array used in the reduction.  Let \(m\) be a
positive integer and identify the \(d=m^2\) coordinates of the frequency
vector with ordered pairs \((i,j)\in[m]\times[m]\).  Let \(B_x\) be a known
polynomial upper bound on every coordinate \(x_{ij}\) in the hard family from
Theorem~\ref{thm:fp-lower}.  Define three pairwise disjoint sets
\[
 A^+:=\{a_1^+,\ldots,a_m^+\},
 \qquad
 A^-:=\{a_1^-,\ldots,a_m^-\},
 \qquad
 C:=\{c_1,\ldots,c_m\},
\]
and put \(A:=A^+\cup A^-\) and \(V:=A\cup C\).  Thus \(|V|=3m\).  Define
\[
 H:=16B_x\left\lceil m^{2/p}\right\rceil.
\]
Choose
\[
 w_x:=\left\lceil\log(H+B_x+1)\right\rceil.
\]
Then every value defined below belongs to \([2^{w_x}-1]\), so the resulting
function is a distance array according to
Section~\ref{sec:preliminaries}.
The distance array \(D_x\) on \(V\) is defined by
\begin{align*}
 D_x(\{u,v\})&:=1
 &&\text{if \(u,v\in A\) or \(u,v\in C\)},\\
 D_x(\{a_i^+,c_j\})&:=H+x_{ij}
 &&\text{if \(i,j\in[m]\)},\\
 D_x(\{a_i^-,c_j\})&:=H-x_{ij}
 &&\text{if \(i,j\in[m]\)}.
\end{align*}

\begin{lemma}[Exact optimum of the reduction]\label{lem:exact-reduction}
For every \(x\in\{0,\ldots,B_x\}^{m\times m}\) and every real \(p\ge1\),
\[
 \OPT_p(D_x)^p=2\Fp(x).
\]
\end{lemma}

\begin{proof}
Define an ultrametric \(U_0\) by
\[
 U_0(\{u,v\}):=
 \begin{cases}
 1,&u,v\in A\text{ or }u,v\in C,\\
 H,&u\in A\text{ and }v\in C.
 \end{cases}
\]
The two sets \(A\) and \(C\) are the two children of the root in a dendrogram
for \(U_0\), so \(U_0\) is an ultrametric.  Its \(p\)-th power cost is
\[
 \sum_{e\in E}|D_x(e)-U_0(e)|^p
 =2\sum_{i,j=1}^m x_{ij}^p
 =2\Fp(x).
\]
In particular,
\[
 \OPT_p(D_x)^p\le2m^2B_x^p.
\]

Let \(U\) be an optimum ultrametric and put \(\tau:=H/2\).  Define a relation
on \(V\) by
\[
 u\sim_U v
 \quad\Longleftrightarrow\quad
 u=v\text{ or }U(\{u,v\})<\tau.
\]
The ultrametric inequality implies that \(\sim_U\) is transitive; hence it is
an equivalence relation.  Suppose that its equivalence classes are not
exactly \(A\) and \(C\).  Then one of the following two events occurs.

First, two vertices in \(A\), or two vertices in \(C\), lie in different
equivalence classes.  For their pair \(e\),
\[
 |D_x(e)-U(e)|\ge\tau-1\ge H/4.
\]
Second, a vertex of \(A\) and a vertex of \(C\) lie in the same equivalence
class.  For their pair \(e\), using \(H\ge4B_x\),
\[
 |D_x(e)-U(e)|
 > H-B_x-\tau
 \ge H/4.
\]
In either case, the \(p\)-th power cost of \(U\) is at least \((H/4)^p\).
The definition of \(H\) gives
\[
 (H/4)^p
 \ge 4^p m^2B_x^p
 >2m^2B_x^p,
\]
contradicting the upper bound supplied by \(U_0\).  Therefore the
equivalence classes of \(\sim_U\) are exactly \(A\) and \(C\).

It follows that all cross-pair values \(U(\{a,c\})\), with \(a\in A\) and
\(c\in C\), are equal.  To verify this statement, let \(a,a'\in A\) and
\(c\in C\).  We have
\[
 U(\{a,a'\})<\tau
 \le\min\{U(\{a,c\}),U(\{a',c\})\}.
\]
If the two cross-pair values were unequal, applying the ultrametric inequality
to the larger one would give a contradiction.  Thus they are equal.  The
same argument with two vertices of \(C\) proves the statement for all cross
pairs.  Denote their common value by \(h\), and write \(h=H+t\).

The contribution of the cross pairs to the \(p\)-th power cost is
\[
 \sum_{i,j=1}^m
 \left(|x_{ij}-t|^p+|x_{ij}+t|^p\right).
\]
For every \(z\ge0\), the function
\[
 g_z(t):=|z-t|^p+|z+t|^p
\]
is convex and even.  Convexity and evenness imply
\[
 g_z(0)
 \le\frac{g_z(t)+g_z(-t)}2
 =g_z(t).
\]
Consequently, the cross-pair contribution is at least
\(2\sum_{i,j}x_{ij}^p=2\Fp(x)\).  The contribution of the remaining pairs is
nonnegative.  Together with the cost of \(U_0\), this proves
\(\OPT_p(D_x)^p=2\Fp(x)\).
\end{proof}

\begin{lemma}[Strict-turnstile realization]\label{lem:strict-realization}
The transformation from an insertion-only stream for \(x\) to a stream for
\(D_x\) is computable online, uses a strict-turnstile stream, and produces
polynomially many updates of polynomial magnitude.
\end{lemma}

\begin{proof}
First insert \(1\) into every coordinate corresponding to a pair contained in
\(A\) or contained in \(C\).  Insert \(H\) into every coordinate
corresponding to a cross pair.  Whenever the frequency stream increments
\(x_{ij}\), add \(1\) to the coordinate \(\{a_i^+,c_j\}\) and add \(-1\) to
the coordinate \(\{a_i^-,c_j\}\).

At every time, \(0\le x_{ij}\le B_x<H\).  Hence every coordinate
\(H-x_{ij}\) remains positive.  All remaining coordinates are also
nonnegative.  The transformed stream is therefore strict turnstile.  It has
\(O(m^2)\) initialization updates and two updates for each update of the
frequency stream.  Since \(B_x\) and the frequency-stream length are
polynomial in \(m\), the transformed stream has polynomial length and
polynomial update magnitude.  Moreover,
\[
 1\le D_x(e)\le H+B_x
\]
for every final coordinate.  Hence every final coordinate has
\(O(\log m)=O(\log n)\) bits.
\end{proof}

\begin{proof}[Proof of Theorem~\ref{thm:lower-intro}]
Fix \(p>4\), and suppose that an algorithm \(\mathcal A\) satisfies the
conclusion required in Theorem~\ref{thm:lower-intro}.  Given an
insertion-only stream for \(x\in\Z_{\ge0}^{m^2}\), use Lemma
\ref{lem:strict-realization} to feed the corresponding strict-turnstile stream
for \(D_x\) to \(\mathcal A\).  Let \(Z\) be the value returned by
\(\mathcal A\), and return
\[
 Y:=\frac{Z^p}{2}
\]
as an estimate of \(\Fp(x)\).

By Lemma~\ref{lem:exact-reduction}, on every successful execution of
\(\mathcal A\),
\[
 (1-\varepsilon_p)^p\Fp(x)
 \le Y\le
 (1+\varepsilon_p)^p\Fp(x).
\]
Choose \(\varepsilon_p>0\) sufficiently small that
\[
 (1-\varepsilon_p)^p\ge1-\gamma_p
 \quad\text{and}\quad
 (1+\varepsilon_p)^p\le1+\gamma_p,
\]
where \(\gamma_p\) is the constant in Theorem~\ref{thm:fp-lower}.  The
resulting procedure is an \(F_p\)-estimation algorithm with success
probability at least \(2/3\).  Theorem~\ref{thm:fp-lower}, with \(d=m^2\),
therefore gives the space lower bound
\[
 \widetilde\Omega\bigl((m^2)^{1-2/p}\bigr)
 =\widetilde\Omega(m^{2-4/p})
 =\widetilde\Omega(n^{2-4/p}),
\]
because \(n=3m\).  Lemma~\ref{lem:strict-realization} proves the asserted
stream restrictions.  The reduction does not restrict the post-processing
time, which completes the proof.
\end{proof}

\section{Concluding statement}

For \(q\in\{0,1,2\}\), a finite canonical family and a high-confidence linear
norm sketch give a one-pass \((1+\varepsilon)\)-approximation in
\(\widetilde O(n/\varepsilon^2)\) bits when the input numbers have
\(O(\log n)\) bits.  For every fixed \(p>4\), reporting a sufficiently accurate
multiplicative estimate of the optimum \(\ell_p\) fitting value in the
strict-turnstile model requires
\(\widetilde\Omega(n^{2-4/p})\) bits.

\small
\begin{thebibliography}{99}

\bibitem{AlonMatiasSzegedy1999}
N.~Alon, Y.~Matias, and M.~Szegedy.
\newblock The space complexity of approximating the frequency moments.
\newblock \emph{Journal of Computer and System Sciences}, 58(1):137--147,
1999.

\bibitem{BarYossefJayramKumarSivakumar2004}
Z.~Bar-Yossef, T.~S. Jayram, R.~Kumar, and D.~Sivakumar.
\newblock An information statistics approach to data stream and communication
complexity.
\newblock \emph{Journal of Computer and System Sciences}, 68(4):702--732,
2004.

\bibitem{CarmelDasKipouridisPipis2025}
A.~Carmel, D.~Das, E.~Kipouridis, and E.~Pipis.
\newblock Fitting tree metrics and ultrametrics in data streams.
\newblock In \emph{52nd International Colloquium on Automata, Languages, and
Programming (ICALP 2025)}, volume 334 of LIPIcs, pages 42:1--42:21, 2025.

\bibitem{CormodeDatarIndykMuthukrishnan2003}
G.~Cormode, M.~Datar, P.~Indyk, and S.~Muthukrishnan.
\newblock Comparing data streams using Hamming norms.
\newblock \emph{IEEE Transactions on Knowledge and Data Engineering},
15(3):529--541, 2003.

\bibitem{Ganguly2011}
S.~Ganguly.
\newblock A lower bound for estimating high moments of a data stream.
\newblock arXiv:1201.0253, 2011.

\bibitem{Indyk2006}
P.~Indyk.
\newblock Stable distributions, pseudorandom generators, embeddings, and data
stream computation.
\newblock \emph{Journal of the ACM}, 53(3):307--323, 2006.

\bibitem{KaneNelsonWoodruff2010Distinct}
D.~M. Kane, J.~Nelson, and D.~P. Woodruff.
\newblock An optimal algorithm for the distinct elements problem.
\newblock In \emph{Proceedings of the 29th ACM SIGMOD-SIGACT-SIGART Symposium
on Principles of Database Systems (PODS 2010)}, pages 41--52, 2010.

\bibitem{KaneNelsonWoodruff2010Norms}
D.~M. Kane, J.~Nelson, and D.~P. Woodruff.
\newblock On the exact space complexity of sketching and streaming small
norms.
\newblock In \emph{Proceedings of the 21st Annual ACM--SIAM Symposium on
Discrete Algorithms (SODA 2010)}, pages 1161--1178, 2010.

\end{thebibliography}

\end{document}
